% Chunk: conclusion of §7.5 and §7.6 Constraints and Dirac Brackets
% Source: Weinberg QFT Vol. I, printed pp. 325--331 (physical PDF pp. 348--354)
% Status: transcribed; source-reviewed; compile-clean; render-reviewed
% Semantic end: paragraph following Eq. (7.6.30), before the §7.7 heading.
% Figures: none. Uncertainties: none.

we must adopt the anticommutation relations
\begin{equation}
\{a_{\mathbf p,\sigma},a^\dagger_{\mathbf p',\sigma'}\}
=\{b_{\mathbf p,\sigma},b^\dagger_{\mathbf p',\sigma'}\}
=\delta^3(\mathbf p'-\mathbf p)\delta_{\sigma'\sigma},
\tag{7.5.49}\label{eq:7.5.49}
\end{equation}
\begin{align}
\{a_{\mathbf p,\sigma},a_{\mathbf p',\sigma'}\}
=\{b_{\mathbf p,\sigma},b_{\mathbf p',\sigma'}\}
&={}\notag\\
\{a_{\mathbf p,\sigma},b_{\mathbf p',\sigma'}\}
=\{a_{\mathbf p,\sigma},b^\dagger_{\mathbf p',\sigma'}\}&=0,
\tag{7.5.50}\label{eq:7.5.50}
\end{align}
and their adjoints. These agree with the results obtained in Chapter 5, thus
verifying that \eqref{eq:7.5.37} is the correct free-particle Hamiltonian for spin
$1/2$. In terms of the $a$s and $b$s, this Hamiltonian is
\begin{equation}
H_0=\sum_\sigma\int \dd^3p\,p^0
\left(a^\dagger_{\mathbf p,\sigma}a_{\mathbf p,\sigma}
-b_{\mathbf p,\sigma}b^\dagger_{\mathbf p,\sigma}\right).
\tag{7.5.51}\label{eq:7.5.51}
\end{equation}
We can rewrite this as a more conventional free-particle Hamiltonian, plus
another infinite c-number\footnote{Note the negative sign of the c-number
term. The conjectured symmetry known as supersymmetry~\hyperref[ref:7.4]{[4]} connects the numbers
of boson and fermion fields, in such a way that the c-numbers in $H_0$ all
cancel.}
\begin{equation}
H_0=\sum_\sigma\int \dd^3p\,p^0\left[
a^\dagger_{\mathbf p,\sigma}a_{\mathbf p,\sigma}
+b^\dagger_{\mathbf p,\sigma}b_{\mathbf p,\sigma}
-\delta^3(\mathbf p-\mathbf p)\right].
\tag{7.5.52}\label{eq:7.5.52}
\end{equation}
The c-number term in Eq.~\eqref{eq:7.5.52} is only important if we worry about
gravitational phenomena; otherwise here, as for the scalar field, we can
throw it away, since it only affects the zero of energy with respect to which
all energies are measured. With this understanding, $H_0$ is a positive
operator, just as for bosons.

\subsection{Constraints and Dirac Brackets}

The chief obstacle to deriving the Hamiltonian from the Lagrangian is the
occurrence of constraints. The standard analysis of this problem is that of
Dirac~\hyperref[ref:7.5]{[5]}, whose terminology we will follow here. Dirac's analysis is not
really needed for the simple theories discussed in this chapter, where it is
easy to identify the unconstrained canonical variables. We shall use the
theory of a real massive vector field for illustration here, returning to
Dirac's approach in the next chapter, where it will be actually useful.

\emph{Primary constraints} are either imposed on the system (as when in the
next chapter we choose a gauge for the electromagnetic field) or arise from
the structure of the Lagrangian itself. For an example of the latter type,
consider the Lagrangian \eqref{eq:7.5.11} of a massive vector field $V^\mu$ interacting
with a current $J_\mu$:
\begin{equation}
\mathcal L=-\frac14F_{\mu\nu}F^{\mu\nu}
-\frac12m^2V_\mu V^\mu-J_\mu V^\mu
\tag{7.6.1}\label{eq:7.6.1}
\end{equation}
where
\begin{equation}
F_{\mu\nu}=\partial_\mu V_\nu-\partial_\nu V_\mu.
\tag{7.6.2}\label{eq:7.6.2}
\end{equation}
Suppose we try to treat all four components of $V^\mu$ on the same basis. We
should then define the conjugates
\begin{equation}
\Pi_\mu\equiv\frac{\partial\mathcal L}{\partial(\partial_0V^\mu)}=-F^{0\mu}.
\tag{7.6.3}\label{eq:7.6.3}
\end{equation}
We immediately find the primary constraint:
\begin{equation}
\Pi_0=0.
\tag{7.6.4}\label{eq:7.6.4}
\end{equation}
More generally, we encounter primary constraints whenever the equations
$\Pi_\ell=\delta L/\delta\partial_0\Psi^\ell$ cannot be solved to give all the
$\partial_0\Psi^\ell$ (at least locally) in terms of $\Pi_\ell$ and
$\Psi^\ell$. This will be the case if and only if the matrix
$\delta^2L/\delta(\partial_0\Psi^\ell)\delta(\partial_0\Psi^m)$ has vanishing
determinant. Such Lagrangians are called \emph{irregular}.

Then there are \emph{secondary constraints}, which arise from the requirement
that the primary constraints be consistent with the equations of motion. For
the massive vector field, this is just the Euler--Lagrange equation~\eqref{eq:7.5.16}
for $V^0$:
\begin{equation}
\partial_i\Pi_i=m^2V^0+J^0.
\tag{7.6.5}\label{eq:7.6.5}
\end{equation}
Here we are finished, but in other theories we might encounter further
constraints by requiring consistency of the secondary constraints with the
field equations, and so on. The distinction between primary, secondary, etc.
constraints is not important; we will treat them all together here.

There is another distinction between certain types of constraint that is more
important. The constraints we have found for the massive vector field are of
a type known as \emph{second class}, for which there is a universal
prescription for the commutation relations. To explain the distinction
between first and second class constraints, and the prescription used to deal
with second class constraints, it is useful first to recall the definition of
the Poisson brackets of classical mechanics.

Consider any Lagrangian $L(\Psi,\dot\Psi)$ that depends on a set of variables
$\Psi^a(t)$ and their time-derivatives $\dot\Psi^a(t)$. (The Lagrangians of
quantum field theory are a special case, with the index $a$ running over all
pairs of $\ell$ and $\mathbf x$.) We can define canonical conjugates for
\emph{all} of these variables by
\begin{equation}
\Pi_a\equiv\frac{\partial L}{\partial\dot\Psi^a}.
\tag{7.6.6}\label{eq:7.6.6}
\end{equation}
The $\Pi$s and $\Psi$s will in general not be independent variables, but may
instead be related by various constraint equations, both primary and
secondary. The Poisson bracket is then defined by
\begin{equation}
[A,B]_{\mathrm P}\equiv
\frac{\partial A}{\partial\Psi^a}\frac{\partial B}{\partial\Pi_a}
-\frac{\partial B}{\partial\Psi^a}\frac{\partial A}{\partial\Pi_a}
\tag{7.6.7}\label{eq:7.6.7}
\end{equation}
with the constraints ignored in calculating the derivatives with respect to
$\Psi^a$ and $\Pi_a$. In particular, we always have
$[\Psi^a,\Pi_b]_{\mathrm P}=\delta^a_b$. (Here and below all fields are taken
at the same time, and time arguments are everywhere dropped.) These brackets
have the same algebraic properties as commutators:
\begin{equation}
[A,B]_{\mathrm P}=-[B,A]_{\mathrm P},
\tag{7.6.8}\label{eq:7.6.8}
\end{equation}
\begin{equation}
[A,BC]_{\mathrm P}=[A,B]_{\mathrm P}C+B[A,C]_{\mathrm P},
\tag{7.6.9}\label{eq:7.6.9}
\end{equation}
including the Jacobi identity
\begin{equation}
[A,[B,C]_{\mathrm P}]_{\mathrm P}
+[B,[C,A]_{\mathrm P}]_{\mathrm P}
+[C,[A,B]_{\mathrm P}]_{\mathrm P}=0.
\tag{7.6.10}\label{eq:7.6.10}
\end{equation}
If we could adopt the usual commutation relations
$[\Psi^a,\Pi_b]=\ii\delta^a_b$, $[\Psi^a,\Psi^b]=[\Pi_a,\Pi_b]=0$, then the
commutator of any two functions of the $\Psi$s and $\Pi$s would be just
$[A,B]=\ii[A,B]_{\mathrm P}$. But the constraints do not always allow this.

The constraints may in general be expressed in the form $\chi_N=0$, where
the $\chi_N$ are a set of functions of the $\Psi$s and $\Pi$s. Because we are
including secondary constraints along with the primary constraints, the set
of all the constraints is necessarily consistent with the equations of
motion $\dot A=[A,H]_{\mathrm P}$, and therefore
\begin{equation}
[\chi_N,H]_{\mathrm P}=0
\tag{7.6.11}\label{eq:7.6.11}
\end{equation}
when the constraint equations $\chi_N=0$ are satisfied.

We call a constraint \emph{first class} if its Poisson bracket with all the
other constraints vanishes when (\emph{after} calculating the Poisson
brackets) we impose the constraints. We shall see a simple example of such a
constraint in the quantization of the electromagnetic field in the next
chapter, where the first class constraint arises from a symmetry of the
action, electromagnetic gauge invariance. In fact, the set of first class
constraints $\chi_N=0$ is always associated with a group of symmetries, under
which an arbitrary quantity $A$ undergoes the infinitesimal transformation
\begin{equation}
\delta_N A\equiv\sum_N\epsilon_N[\chi_N,A]_{\mathrm P}.
\tag{7.6.12}\label{eq:7.6.12}
\end{equation}
(In field theory these are local transformations, because the index $N$
contains a spacetime coordinate.) Eq.~\eqref{eq:7.6.11} shows that this transformation
leaves the Hamiltonian invariant, and for first class constraints it also
respects all other constraints. Such first class constraints can be
eliminated by a choice of gauge, or treated by gauge-invariant methods
described in Volume II.

After all of the first class constraints have been eliminated, the remaining
constraint equations $\chi_N=0$ are such that no linear combination
$\sum_Nu_N[\chi_N,\chi_M]_{\mathrm P}$ of the Poisson brackets of these
constraints with each other vanishes. It follows that the matrix of the
Poisson brackets of the remaining constraints is non-singular:
\begin{equation}
\operatorname{Det}C\ne0,
\tag{7.6.13}\label{eq:7.6.13}
\end{equation}
where
\begin{equation}
C_{NM}\equiv[\chi_N,\chi_M]_{\mathrm P}.
\tag{7.6.14}\label{eq:7.6.14}
\end{equation}
Constraints of this sort are called \emph{second class}. Note that there must
always be an even number of second class constraints, because an
antisymmetric matrix of odd dimensionality necessarily has vanishing
determinant.

As we have seen, in the case of the massive real vector field the constraints
are
\begin{equation}
\chi_{1\mathbf x}=\chi_{2\mathbf x}=0,
\tag{7.6.15}\label{eq:7.6.15}
\end{equation}
where
\begin{equation}
\chi_{1\mathbf x}=\Pi_0(\mathbf x),\qquad
\chi_{2\mathbf x}=\partial_i\Pi_i(\mathbf x)-m^2V^0(\mathbf x)-J^0(\mathbf x).
\tag{7.6.16}\label{eq:7.6.16}
\end{equation}
The Poisson bracket of these constraints is
\begin{equation}
C_{1\mathbf x,2\mathbf y}=-C_{2\mathbf y,1\mathbf x}
=[\chi_{1\mathbf x},\chi_{2\mathbf y}]_{\mathrm P}
=m^2\delta^3(\mathbf x-\mathbf y)
\tag{7.6.17}\label{eq:7.6.17}
\end{equation}
and, of course,
\begin{equation}
C_{1\mathbf x,1\mathbf y}=C_{2\mathbf x,2\mathbf y}=0.
\tag{7.6.18}\label{eq:7.6.18}
\end{equation}
This `matrix' is obviously non-singular, so the constraints \eqref{eq:7.6.15} are
second class.

Dirac suggested that when all constraints are second class, the commutation
relations will be given by
\begin{equation}
[A,B]=\ii[A,B]_{\mathrm D},
\tag{7.6.19}\label{eq:7.6.19}
\end{equation}
where $[A,B]_{\mathrm D}$ is a generalization of the Poisson bracket known as
the \emph{Dirac bracket}:
\begin{equation}
[A,B]_{\mathrm D}\equiv[A,B]_{\mathrm P}
-[A,\chi_N]_{\mathrm P}(C^{-1})^{NM}[\chi_M,B]_{\mathrm P}.
\tag{7.6.20}\label{eq:7.6.20}
\end{equation}
(Here $N$ and $M$ are compound indices including the position in space,
taking values like $1,\mathbf x$ and $2,\mathbf x$ in the vector field
example.) He noted that the Dirac bracket like the Poisson bracket satisfies
the same algebraic relations as the commutators
\begin{equation}
[A,B]_{\mathrm D}=-[B,A]_{\mathrm D},
\tag{7.6.21}\label{eq:7.6.21}
\end{equation}
\begin{equation}
[A,BC]_{\mathrm D}=[A,B]_{\mathrm D}C+B[A,C]_{\mathrm D},
\tag{7.6.22}\label{eq:7.6.22}
\end{equation}
\begin{equation}
[A,[B,C]_{\mathrm D}]_{\mathrm D}
+[B,[C,A]_{\mathrm D}]_{\mathrm D}
+[C,[A,B]_{\mathrm D}]_{\mathrm D}=0,
\tag{7.6.23}\label{eq:7.6.23}
\end{equation}
and also the relations
\begin{equation}
[\chi_N,B]_{\mathrm D}=0
\tag{7.6.24}\label{eq:7.6.24}
\end{equation}
which make the commutation relations~\eqref{eq:7.6.19} consistent with the constraints
$\chi_N=0$. Also, the Dirac brackets are unchanged if we replace the $\chi_N$
with any functions $\chi'_N$ for which the equations $\chi'_N=0$ and
$\chi_N=0$ define the same submanifold of phase space. But all these
agreeable properties do not prove that the commutators are actually given by
Eq.~\eqref{eq:7.6.19} in terms of the Dirac brackets.

This issue is illuminated if not settled by a powerful theorem proved by
Maskawa and Nakajima~\hyperref[ref:7.6]{[6]}. They showed that for any set of canonical variables
$\Psi^a$, $\Pi_a$ governed by second class constraints, it is always possible
by a canonical transformation\footnote{Recall that by a canonical
transformation, we mean a transformation from a set of phase space coordinates
$\Psi^a$, $\Pi_a$ to some other phase space coordinates
$\widetilde\Psi^a$, $\widetilde\Pi_a$, such that
$[\widetilde\Psi^a,\widetilde\Pi_b]_{\mathrm P}=\delta^a_b$ and
$[\widetilde\Psi^a,\widetilde\Psi^b]_{\mathrm P}
=[\widetilde\Pi_a,\widetilde\Pi_b]_{\mathrm P}=0$, the Poisson brackets being
calculated in terms of the $\Psi^a$ and $\Pi_a$. It follows that the Poisson
brackets for any functions $A,B$ are the same whether calculated in terms of
$\Psi^a$ and $\Pi_a$ or in terms of $\widetilde\Psi^a$ and
$\widetilde\Pi_a$. It also follows that if $\Psi^a$ and $\Pi_a$ satisfy the
Hamiltonian equations of motion, then so do $\widetilde\Psi^a$ and
$\widetilde\Pi_a$, with the same Hamiltonian. The Lagrangian is changed by a
canonical transformation, but only by a time-derivative, which does not affect
the action.} to construct two sets of variables $Q^n$, $\mathcal Q^r$ and
their respective conjugates $P_n$, $\mathcal P_r$, such that the constraints
read $\mathcal Q^r=\mathcal P_r=0$. Using these coordinates to calculate
Poisson brackets, and redefining the constraint functions as
$\chi_{1r}=\mathcal Q^r$, $\chi_{2r}=\mathcal P_r$, we have
\[
C_{1r,2s}=[\mathcal Q^r,\mathcal P_s]_{\mathrm P}=\delta^r_s,
\]
\[
C_{1r,1s}=[\mathcal Q^r,\mathcal Q^s]_{\mathrm P}=0,
\qquad
C_{2r,2s}=[\mathcal P_r,\mathcal P_s]_{\mathrm P}=0,
\]
and for any functions $A,B$
\[
[A,\chi_{1r}]_{\mathrm P}=-\frac{\partial A}{\partial\mathcal P_r},
\qquad
[A,\chi_{2r}]_{\mathrm P}=\frac{\partial A}{\partial\mathcal Q^r}.
\]
This $C$-matrix has inverse $C^{-1}=-C$, so the Dirac brackets \eqref{eq:7.6.20} are
here
\begin{align}
[A,B]_{\mathrm D}
&=[A,B]_{\mathrm P}
+[A,\chi_{1r}]_{\mathrm P}[\chi_{2r},B]_{\mathrm P}
-[A,\chi_{2r}]_{\mathrm P}[\chi_{1r},B]_{\mathrm P}\notag\\
&=[A,B]_{\mathrm P}
-\frac{\partial A}{\partial\mathcal Q^r}
 \frac{\partial B}{\partial\mathcal P_r}
+\frac{\partial B}{\partial\mathcal Q^r}
 \frac{\partial A}{\partial\mathcal P_r}\notag\\
&=\frac{\partial A}{\partial Q^n}\frac{\partial B}{\partial P_n}
-\frac{\partial B}{\partial Q^n}\frac{\partial A}{\partial P_n}.
\tag{7.6.25}\label{eq:7.6.25}
\end{align}
In other words, \emph{the Dirac bracket is equal to the Poisson bracket
calculated in terms of the reduced set of unconstrained canonical variables
$Q^n$, $P_n$}. If we assume that these unconstrained variables satisfy the
canonical commutation relations, then the commutators of general operators
$A,B$ are given by Eq.~\eqref{eq:7.6.19} in terms of the Dirac brackets.\footnote{It is
still an open question whether we should adopt canonical commutation relations
for the unconstrained variables $Q^n$, $P_n$ constructed by the
Maskawa--Nakajima canonical transformation. Ultimately, the test of such
canonical commutation relations is their consistency with the free-field
commutation relations derived in Chapter 5, but to apply this test we need to
know what the $Q^n$ and $P_n$ are. In the Appendix to this chapter we display
two large classes of theories in which we can identify a set of unconstrained
$Q$s and $P$s, such that the Dirac commutation relations~\eqref{eq:7.6.19} follow from
the ordinary canonical commutation relations of the $Q$s and $P$s. We shall
also show that in these cases, the Hamiltonian defined in terms of the
unconstrained $\Psi$s and $\Pi$s may be written just as well in terms of the
constrained variables.}

We now return to the massive vector field, to see how it can be quantized
using Dirac brackets. This is a case where it is easy to express the
constrained variables $V^0$ and $\Pi_0$ in terms of the unconstrained
ones\footnote{This is a special case of the theories discussed in Part A of
the Appendix.} $V_i$ and $\Pi_i$; we have simply $\Pi_0=0$, and $V^0$ is
given by Eq.~\eqref{eq:7.6.5}. From Eqs.~\eqref{eq:7.6.17} and \eqref{eq:7.6.18}, we see that $C_{NM}$
here has the inverse
\begin{equation}
(C^{-1})^{1\mathbf x,2\mathbf y}
=-(C^{-1})^{2\mathbf y,1\mathbf x}
=-m^{-2}\delta^3(\mathbf x-\mathbf y),
\tag{7.6.26}\label{eq:7.6.26}
\end{equation}
\begin{equation}
(C^{-1})^{1\mathbf x,1\mathbf y}
=(C^{-1})^{2\mathbf x,2\mathbf y}=0.
\tag{7.6.27}\label{eq:7.6.27}
\end{equation}
Therefore the Dirac prescription \eqref{eq:7.6.19}, \eqref{eq:7.6.20} yields the equal-time
commutators
\begin{align}
[A,B]={}&\ii[A,B]_{\mathrm P}\notag\\
&+\ii m^{-2}\int \dd^3z\Bigl(
[A,\Pi_0(\mathbf z)]_{\mathrm P}
[\partial_i\Pi_i(\mathbf z)-m^2V^0(\mathbf z)-J^0(\mathbf z),B]_{\mathrm P}
-A\leftrightarrow B\Bigr).
\tag{7.6.28}\label{eq:7.6.28}
\end{align}
By definition, we have
\begin{equation}
[V^\mu(\mathbf x),\Pi_\nu(\mathbf y)]_{\mathrm P}
=\delta^3(\mathbf x-\mathbf y)\delta^\mu_\nu,
\qquad
[V^\mu(\mathbf x),V^\nu(\mathbf y)]_{\mathrm P}
=[\Pi_\mu(\mathbf x),\Pi_\nu(\mathbf y)]_{\mathrm P}=0.
\tag{7.6.29}\label{eq:7.6.29}
\end{equation}
Hence
\begin{align}
[V^i(\mathbf x),V^j(\mathbf y)]&=[V^0(\mathbf x),V^0(\mathbf y)]=0,
\notag\\
[V^i(\mathbf x),V^0(\mathbf y)]&=-\ii m^{-2}\partial_i\delta^3(\mathbf x-\mathbf y),
\notag\\
[V^i(\mathbf x),\Pi_j(\mathbf y)]&=\ii\delta^i_j\delta^3(\mathbf x-\mathbf y),
\tag{7.6.30}\label{eq:7.6.30}\\
[V^0(\mathbf x),\Pi_j(\mathbf y)]&=[V^\mu(\mathbf x),\Pi_0(\mathbf y)]=0,
\notag\\
[\Pi^\mu(\mathbf x),\Pi^\nu(\mathbf y)]&=0.
\notag
\end{align}
These are indeed just the commutation relations that we would find by
assuming that the unconstrained variables satisfy the usual canonical
commutation relations
$[V^i(\mathbf x),\Pi_j(\mathbf y)]=\ii\delta^i_j\delta^3(\mathbf x-\mathbf y)$,
$[V^i(\mathbf x),V^j(\mathbf y)]=[\Pi_i(\mathbf x),\Pi_j(\mathbf y)]=0$,
and using the constraints to evaluate the commutators involving $\Pi_0$ and
$V^0$.
